\begin{derivation}[从\eqnrefs{eq:sm-yukawa-general-dp11}到\eqnrefs{eq:higgs-fermion-coupling-dp11}]
对任意一种带电费米子 $f=u,d,e$，电弱破缺后正文 Yukawa 项都具有同一矩阵结构
\begin{equation*}
 \Lag_{Y,f}
 =-\frac{v_E}{\sqrt2}
 \left(\bar f_LY_ff_R+\bar f_RY_f^\dagger f_L\right)
 -\frac{\sqrt{\hbar c}\,h}{\sqrt2}
 \left(\bar f_LY_ff_R+\bar f_RY_f^\dagger f_L\right).
\end{equation*}
这里 $Y_f$ 是世代空间中的任意复 $N_g\times N_g$ 矩阵；它一般既不厄米也不正规，因此不能假定存在单个幺正矩阵 $U$ 使 $U^\dagger Y_fU$ 对角。

从两个厄米半正定矩阵
\begin{equation*}
 H_L\equiv Y_fY_f^\dagger,
 \qquad
 H_R\equiv Y_f^\dagger Y_f
\end{equation*}
出发。谱定理保证存在幺正矩阵 $U_{fL}$ 使
\begin{equation*}
 U_{fL}^\dagger H_LU_{fL}
 =\operatorname{diag}(y_1^2,\ldots,y_{N_g}^2),
 \qquad y_i\ge0.
\end{equation*}
设 $u_{Li}$ 是 $U_{fL}$ 的第 $i$ 个列向量，即
\begin{equation*}
 H_Lu_{Li}=y_i^2u_{Li}.
\end{equation*}
对任意 $y_i>0$，定义右奇异向量
\begin{equation*}
 u_{Ri}\equiv\frac1{y_i}Y_f^\dagger u_{Li}.
\end{equation*}
它们的内积为
\begin{align*}
 u_{Ri}^\dagger u_{Rj}
 &=\frac1{y_iy_j}
 u_{Li}^\dagger Y_fY_f^\dagger u_{Lj}\\
 &=\frac{y_j^2}{y_iy_j}u_{Li}^\dagger u_{Lj}
 =\delta_{ij},
\end{align*}
并且
\begin{align*}
 Y_fu_{Ri}
 &=\frac1{y_i}Y_fY_f^\dagger u_{Li}
 =y_i u_{Li},\\
 Y_f^\dagger u_{Li}
 &=y_i u_{Ri}.
\end{align*}
若存在 $y_i=0$，在 $\ker Y_f$ 与 $\ker Y_f^\dagger$ 中分别补齐正交基即可。把全部 $u_{Ri}$ 作为列向量组成幺正矩阵 $U_{fR}$，便有逐矩阵元关系
\begin{align*}
 (U_{fL}^\dagger Y_fU_{fR})_{ij}
 &=u_{Li}^\dagger Y_fu_{Rj}\\
 &=y_j u_{Li}^\dagger u_{Lj}
 =y_j\delta_{ij}.
\end{align*}
因此
\begin{equation*}
 \boxed{
 U_{fL}^\dagger Y_fU_{fR}
 =\operatorname{diag}(y_1,\ldots,y_{N_g})}
\end{equation*}
就是正文\eqnrefs{eq:yukawa-svd-dp11}。奇异值按定义非负，因此由 $m_f=v_E y_f/(\sqrt2 c)$ 得到的质量本征值可统一取非负实数。

定义新的左右手场
\begin{equation*}
 f_L'=U_{fL}^\dagger f_L,
 \qquad
 f_R'=U_{fR}^\dagger f_R.
\end{equation*}
因为 $U_{fL},U_{fR}$ 只作用在世代空间而且为常数矩阵，动能保持不变：
\begin{align*}
 \bar f_L\ii\hbar c\gamma^\mu D_\mu f_L
 &=\bar f_L'U_{fL}^\dagger
 \ii\hbar c\gamma^\mu D_\mu U_{fL}f_L'\\
 &=\bar f_L'\ii\hbar c\gamma^\mu D_\mu f_L',
\end{align*}
右手项同理。这里使用的是规范相互作用在同一电荷扇区的世代空间中正比于单位矩阵，因此 $D_\mu$ 与 $U_{fL,R}$ 对易。

质量项变为
\begin{align*}
 -\frac{v_E}{\sqrt2}\bar f_LY_ff_R+\mathrm{H.c.}
 &=-\frac{v_E}{\sqrt2}
 \bar f_L'U_{fL}^\dagger Y_fU_{fR}f_R'
 +\mathrm{H.c.}\\
 &=-\sum_i\frac{y_iv_E}{\sqrt2}
 \left(\bar f_{Li}'f_{Ri}'+\bar f_{Ri}'f_{Li}'\right)\\
 &=-\sum_i m_{fi}c^2\,\bar f_i'f_i',
\end{align*}
其中
\begin{equation*}
 \boxed{m_{fi}c^2=\frac{y_iv_E}{\sqrt2}},
\end{equation*}
即正文\eqnrefs{eq:fermion-masses-yukawa-dp11}。

关键点是 Higgs 涨落项含有\emph{完全相同的}矩阵 $Y_f$，因此不需要第二次对角化：
\begin{align*}
 \Lag_{hff}
 &=-\frac{\sqrt{\hbar c}\,h}{\sqrt2}
 \bar f_L'U_{fL}^\dagger Y_fU_{fR}f_R'
 +\mathrm{H.c.}\\
 &=-\sum_i
 \frac{y_i\sqrt{\hbar c}}{\sqrt2}
 h\left(\bar f_{Li}'f_{Ri}'+\bar f_{Ri}'f_{Li}'\right)\\
 &=-\sum_i
 \frac{m_{fi}c^2}{v_E}\sqrt{\hbar c}\,
 h\,\bar f_i'f_i'.
\end{align*}
这就是正文\eqnrefs{eq:higgs-fermion-coupling-dp11}。所以树级 Higgs 中性耦合在质量基中自动对角，且耦合强度与费米子质量线性成正比。

若某组奇异值严格简并，则对应简并子空间内仍允许额外幺正旋转而不改变质量矩阵；因此 CKM/PMNS 的独立参数计数需要单独注明“质量谱非简并”这一条件。质量本征值的定义和 Higgs 中性耦合的对角性本身并不依赖这一假设。
\end{derivation}
